
\documentclass{article} % For LaTeX2e
\usepackage[submission]{colm2026_conference}

\usepackage{microtype}
\usepackage{hyperref}
\usepackage{url}
\usepackage{booktabs}
\usepackage{enumitem}
\usepackage{tabularx}
\usepackage[most]{tcolorbox}
\usepackage{graphicx}
\usepackage{xcolor}
\usepackage{listings}
\usepackage{algorithm}
\usepackage{algorithmic}
\usepackage{amsmath,amsfonts,amssymb}
\usepackage{subcaption}
\usepackage{tikz}
\usetikzlibrary{shapes.geometric, arrows.meta, positioning, calc, fit}
\usepackage{caption}

\usepackage{lineno}

\definecolor{darkblue}{rgb}{0, 0, 0.5}
\hypersetup{colorlinks=true, citecolor=darkblue, linkcolor=darkblue, urlcolor=darkblue}

% Custom commands
\newcommand{\K}{\mathcal{K}}
\newcommand{\agents}{\mathcal{A}}
\newcommand{\objects}{\mathcal{O}}
\newcommand{\rooms}{\mathcal{R}}
\newcommand{\mechanics}{\mathcal{M}}
\newcommand{\task}{\mathcal{T}}
\newcommand{\goal}{\varphi}
\newcommand{\taskdesc}{\delta}
\newcommand{\pddl}{\textsc{pddl}}
\newcommand{\saaket}[1]{\textcolor{blue}{[Saaket: #1]}}
\newcommand{\dylan}[1]{\textcolor{blue}{[Dylan: #1]}}
\lstset{
  basicstyle=\ttfamily\small,
  breaklines=true,
  frame=single,
  columns=fullflexible,
  keepspaces=true,
  xleftmargin=1em,
  framexleftmargin=0.5em
}

\title{EmToM: An Evolving Benchmark for Embodied Theory of Mind}

\author{}

\newcommand{\fix}{\marginpar{FIX}}
\newcommand{\new}{\marginpar{NEW}}
\newcommand{\g}[1]{\textcolor{red}{#1}}
\newcommand{\shortname}{\texttt{EmToM}}
\begin{document}

\ifcolmsubmission
\linenumbers
\fi

\maketitle

\begin{abstract}
Theory of Mind, the ability to track others' epistemic state, is what makes humans efficient decentralized collaborators. This ability is also required by AI agents in multi-agent settings to avoid duplicating effort, communicate efficiently, anticipate what other agents will do and plan accordingly. 
Existing benchmarks measure literal Theory of Mind by prompting agents with explicit questions about other's beliefs, but fail to measure functional Theory of Mind, whether agents can infer and act on others' mental states without being asked.  
We present \shortname{}, an benchmark to measure functional Theory of Mind in embodies agents. These agents operate in 3D scenes under partial observability and information asymmettery.  Each task in \shortname{} is formally verified to be solvable and to require a specific depth of epistemic reasoning; further, new tasks are generated to increase difficulty as models  
improve. We evaluate frontier agents backed by the GPT and Claude families on over 500 \shortname{} tasks and find that these agents have a pass rate less than 25\%, with performance dropping sharply at second-order epistemic depth. Qualitative analysis reveals that agents fail to propagate beliefs through required knowledge chains, and do not maintain accurate models of what their partners can observe or access, leading to misdirected delegation and absence of strategic reasoning.                                                       
   


% Evaluating GPT-5.2 and Sonnet on 68 tasks, both pass under 25\%: agents exhaust limited messages on coordination overhead instead of task-critical observations, and spiral into recursive
% confirmation loops where each agent asks the next to verify before anyone acts. These patterns — along with zero strategic defection and purely reactive competitive play — expose a gap between fluent language and grounded belief reasoning that warrants further investigation.
\end{abstract}

\section{Introduction}
\label{sec:introduction}
\g{questions to answer - (a) do we really need ToM in agents? Can't we do by good prompting? How does it make our lives easier if agents have ToM? Why are we simulating constraints like limited communication bandwidth?(b) what would count as a ToM failure and what as a reasoning failure (c) Types of ToM grounded in literature? }
Humans build internal representations of other people's beliefs, intentions, and mental states, and use them to predict behavior. This capacity is called Theory of Mind (ToM)~\citep{premack1978does, wimmer1983beliefs}. By age four, children recognize that others can hold false beliefs. For instance, they grasp that someone who left the room still believes an object is where they last saw it~\citep{wimmer1983beliefs, tomasello2005understanding}. Adults rely on ToM constantly for coordination. When two people carry a couch through a doorway, each predicts how the other will move based on what the other can see. ToM makes efficient coordination possible under partial observability and information asymmetry~\citep{tomasello2005understanding, bratman1992shared}. \saaket{Do we need to specifically mention coordination here? It will set the tone of the benchmark. How about something like: Adults rely on ToM constantly in everyday interactions, adjusting what they say and do based on what they believe others know, want, or expect}

Multi-agent systems deployed in embodied environments, such as household robots coordinating to clean a room or search-and-rescue teams navigating a building, require the same capacity. An agent that cannot model what its partner knows will broadcast information the partner already has, fail to anticipate partner actions, and waste time waiting for instructions that never arrive. 
Consider two agents searching a house for an object: without ToM, both search the same rooms; with ToM, each agent infers what the other has already checked and covers new ground. 
Embodied collaboration under partial observability requires agents to reason about each other's epistemic states, i.e. what has each agent seen, what can each agent infer, and what must be communicated. This is called functional ToM, active belief modeling that drives action selection during coordination. \saaket{sentences feel very loaded, too much to keep in working memory when reading. I think these two paragraphs are kind of repeating. We can trim this down. }

Existing benchmarks fail to test functional ToM, especially in embodied settings. Text-based evaluations like Sally-Anne variants~\citep{wimmer1983beliefs}, ToMi~\citep{le2019revisiting}, HiToM~\citep{wu2023hi}, BigToM~\citep{gandhi2024understanding}, FANToM~\citep{kim2023fantom}, and ExploreToM~\citep{exploretom2024} test what \citet{riemer2024broken} call \emph{literal} ToM. The model reads a vignette and answers questions about characters' beliefs, without ever acting on those inferences. Multimodal extensions like MMToM-QA~\citep{mmtomqa2024} and MuMA-ToM~\citep{mumatom2024} add video input but remain observational. Embodied benchmarks like PARTNR~\citep{partnr2024} are embodied but evaluate multiagent instruction following rather than epistemic state reasoning. On the other hand, large scale simulations like SOTOPIA~\citep{sotopia2023} tests social intelligence through dialogue but has no embodiment or spatial grounding. \saaket{The specifics should go in related work, we should keep the part that describes what is missing as an objective problem without talking about the benchmarks and papers here.}
All these benchmarks are hand-crafted and fixed-size; as models improve, scores saturate without reflecting genuine advances in ToM understanding.

In this work, we propose EmToM, an \emph{evolving} benchmark for functional ToM in embodied multi-agent settings. EmToM is embodied, having 2-4 agents act in 3D household scenes from Habitat~\citep{szot2021habitat} with egocentric RGB observations under partial observability. It is agentic, requiring the agents to take actions while coordinating to achieve shared, competitive, or mixed-motive goals. It is formally grounded, each task's goal is a \pddl{} formula with epistemic operators ($\K$-nesting up to depth 3), and a verifier confirms both that the task is solvable and that solving it genuinely requires the specified ToM depth. And it evolve, a closed-loop pipeline generates new tasks, verifies them against the simulator, and calibrates difficulty according to the model making the benchmark difficult to saturate without genuine ToM advances. \saaket{Need to clarify why we need evolving benchmarks before mentioning them.}

% TODO: update numbers once full evaluation is done
\g{do it once we have proper results}
We evaluate frontier models (GPT-5.2 with vision, Sonnet with text) on 68 EmToM tasks spanning cooperative, competitive, and mixed-motive settings at ToM depths 1--3. We find that: (i)~second-order ToM is achievable but fragile, succeeding only in simple topologies with generous communication budgets; (ii)~agents exhibit a striking cooperative bias, with zero defection across all mixed-motive tasks; (iii)~communication bandwidth is the primary bottleneck, with 60\%+ of depth-2 failures traced to message exhaustion; and (iv)~agents fall into novel failure modes such as confirmation spirals and action repetition loops that have not been documented in prior ToM evaluations.


%%% Agent sees seed tasks --> writes the pddl --> k depth and task standardization --> write a story (task description) and agent secrets  (Tasks get reused) --> ed has ideas about making seed tasks. 



% Creation
%%% 
%%% Mechanics: Manually created -- attribution to games where we see this. 
%%% t_i = f(Prompt, Mechanics, S(t_0, t_1.. t_{i-1}), Action Space, Scene Objects and Rooms 7(Graph))  
% S is selection which control evolution 
%%% Initial Prompt --> JSON File (includes title, category, task_desc, scenes, mechanics, secrets, action available, PDDL) 
%%% 

% Judging 
%%% PDDL -> [Solver] -> Golden Trajectory, ToM Depth 
%%% Task JSON, Golden Trajectory, Regenerates Golden Traj (Stats, %) -> Judge LLM -> Keep/Don't Keep  (t_i) 
%%% Evaluation(Target Model) 
%%%

%%% {(t_i, solvable?), \phi}   

%%% Take a set of tasks which are solvable 

% ToM --> Baseline vs Real (Make Decision: Should this be included in pipeline or as separate analysis?) 

\section{Preliminaries}
\label{sec:preliminaries}

\subsection{Functional and Literal Theory of Mind}
\label{sec:func_tom}

Literal Theory of Mind evaluations probe an agent's ability to report the beliefs, desires, or intentions of others when asked directly. Functional Theory of Mind, by contrast, requires that agents \emph{act} on their understanding of others' mental states in order to succeed at a task~\citep{street2024llms}. In our setting, agents operate in a shared 3D environment with private information, limited communication, and physical constraints. Success depends not on answering questions about what others believe, but on choosing actions that account for what others know, can observe, and intend to do.


\subsection{Epistemic operators and ToM depth}
\label{sec:epistemic_ops}

Consider a task in which Agent~A must place a bowl on a table, and Agent~B must \emph{confirm} that the bowl has been placed. Agent~B can only confirm what it has observed or been told, so the task goal is not just that the bowl is on the table, but that B \emph{knows} it is. Now suppose a third party, Agent~C, needs to verify that B has confirmed the placement. C must reason about B's knowledge state, which in turn depends on what B observed about A's actions. Each additional layer of ``who knows what about whom'' demands a deeper level of Theory of Mind.

We formalize these knowledge requirements with the epistemic operator $\K$. Writing $\K_i(\phi)$ asserts that agent $a_i$ knows fact $\phi$ by the end of an episode, which is satisfied only when $a_i$ has directly observed $\phi$ or received it through communication. In the example above, the goal that B knows the bowl is placed is written $\K_B(\texttt{is\_on\_top}(\texttt{bowl}, \texttt{table}))$. The further requirement that C knows B knows is written $\K_C(\K_B(\texttt{is\_on\_top}(\texttt{bowl}, \texttt{table})))$. \dylan{We do not check K-goals this way anymore since it is brittle (communication parsing doesn't account for lying or vague natural language). Instead, we directly probe agents by questioning at the end of episode (e.g. what state does agent 0 think the door is in?). K-goals no longer contribute to task success, and are instead "probes" at the end of the episode to test literal ToM.}

These operators nest to arbitrary depth $d$:
%
\begin{equation}
  \underbrace{\K_{a_1}(\K_{a_2}(\cdots \K_{a_d}(\phi) \cdots))}_{\text{depth } d}
  \label{eq:tom_depth}
\end{equation}
%
The nesting depth directly determines the order of ToM reasoning required. At depth~1, an agent simply needs to learn a fact, which requires no modeling of others (zeroth-order). At depth~2, as in the $\K_C(\K_B(\phi))$ example, an agent must reason about another agent's knowledge, requiring first-order ToM. Depth~3 demands reasoning about what one agent believes another agent knows about a third, and so on. In general, depth $d$ corresponds to $(d{-}1)$-th order ToM.

We connect this to the level-$k$ reasoning framework from behavioral game theory~\citep{stahl1994experimental, camerer2004cognitive}. A level-0 agent acts on a fixed heuristic without considering others. A level-$k$ agent chooses actions assuming the other agents reason at level $k{-}1$:
%
\begin{equation}
  \pi^{(0)}_i = \text{fixed heuristic}, \qquad \pi^{(k)}_i = \arg\max_{\pi_i} \; U_i\!\left(\pi_i,\; \pi^{(k-1)}_{-i}\right) \;\; \text{for } k \geq 1
  \label{eq:level_k}
\end{equation}
%

\subsection{Epistemic planning via \pddl{}}
\label{sec:epistemic_pddl}

We use \pddl{} (Planning Domain Definition Language) to formally specify task goals and verify that they are solvable. A goal mixes physical predicates with epistemic $\K$-operators. Consider the goal in Figure~\ref{fig:pddl_example}: the bowl must end up on the table, the cabinet must be open, and agent\_0 must know that agent\_1 knows the bowl has been placed. The physical parts are standard \pddl{}, but the $\K$-goals require reasoning about knowledge and communication, which classical planners do not natively support.

\begin{figure}[h]
\centering
\begin{lstlisting}
(and
  (is_on_top bowl_1 table_22)
  (K agent_0 (K agent_1
    (is_on_top bowl_1 table_22)))
  (is_open cabinet_34)
)
\end{lstlisting}
\caption{Example \pddl{} goal with epistemic nesting ($\K$-depth 2, first-order ToM). The physical subgoals require placing the bowl and opening the cabinet. The epistemic subgoal requires agent\_0 to know that agent\_1 knows the bowl has been placed.}
\label{fig:pddl_example}
\end{figure}

To bridge this gap, we compile epistemic goals into purely classical \pddl{} that a standard planner can solve~\citep{muise2015planning}. No agents are actually running at this stage. The compiler rewrites the epistemic goal into an equivalent classical planning problem, and a planner then searches for an abstract action sequence that satisfies it. If such a sequence exists, the task is provably solvable; if not, the task is rejected.

We illustrate with the goal from Figure~\ref{fig:pddl_example}. The physical subgoal, placing the bowl on the table, is already expressible in classical \pddl{}. The epistemic subgoal, $\K_0(\K_1(\texttt{is\_on\_top}(\texttt{bowl}, \texttt{table})))$, is not. The compiler translates it as follows:

\begin{enumerate}
  \item For each fact that an agent may need to know, the compiler creates a new predicate to represent that knowledge. In this case it creates \texttt{(knows\_agent\_1\_is\_on\_top bowl table)} to represent Agent~1 knowing the bowl is on the table, and \texttt{(knows\_agent\_0\_that\_agent\_1\_knows ...)} to represent Agent~0 knowing that Agent~1 has this knowledge.
  \item The compiler generates abstract planning operators that can set these predicates to true. An \emph{observe} operator, available when an agent is co-located with the relevant objects, sets the first-layer knowledge predicate. An \emph{inform} operator, available when one agent can communicate with another, propagates knowledge between agents and sets the higher-layer predicate. Each inform operator consumes one token from the sending agent's message budget $b_i$, so the planner must respect bandwidth limits.
  \item The original epistemic goal is replaced by a conjunction of these classical predicates. The planner now searches over the physical actions (pick up, place, navigate) together with the observe and inform operators to find a plan that achieves both the physical and knowledge subgoals.
\end{enumerate}

The same pattern extends to deeper nesting: a depth-3 goal adds a third layer of predicates and additional inform operators. The compiled problem is solved by Fast Downward in a single call. The $\K$-depth is computed directly from the goal structure during compilation.




\subsection{Mechanics}
\label{sec:mechanics}

Mechanics are predefined constraint types that can be composed to create information asymmetry and force epistemic reasoning. Each task selects and configures a subset from the following registry:
\begin{itemize}
  \item \textbf{Room restriction} confines agents to subsets of rooms, creating observation asymmetry.
  \item \textbf{Limited bandwidth} caps each agent's messages to $b_i$, forcing communication prioritization.
  \item \textbf{Restricted communication} constrains who can message whom via a directed graph.
  \item \textbf{Remote control} links a trigger object to a target across rooms.
  \item \textbf{State mirroring} synchronizes two objects' states across locations.
  \item \textbf{Inverse state} reverses an object's affordance (opening it closes it).
\end{itemize}
These compose freely. A task may combine room restrictions with limited bandwidth and restricted communication to create multi-hop relay problems where the choice of relay partner is itself the ToM challenge.


\section{Methodology}
\label{sec:methodology}

Evaluating functional Theory of Mind requires tasks where success depends on agents adapting their actions to the private knowledge, access, and intentions of other agents, not merely reporting beliefs when prompted. Hand-crafting such tasks does not scale, and as models improve, fixed benchmarks saturate. We propose an agentic task generation framework that addresses both problems. An autonomous coding agent authors multi-agent ToM tasks inside a sandboxed workspace, invoking verification tools that ensure each task is logically solvable, physically executable, and genuinely requires epistemic reasoning.

We first define the task representation (Section~\ref{sec:task_repr}), then describe the generation agent (Section~\ref{sec:gen_agent}) and its verification tools: quality judging (Section~\ref{sec:judge}), structural validation (Section~\ref{sec:verification}), and empirical calibration (Section~\ref{sec:tom_necessity}).

\subsection{Task representation}
\label{sec:task_repr}

Each task $\task$ is a tuple:
%
\begin{equation}
  \task = \bigl(\mathcal{S},\; \agents,\; \goal,\; \taskdesc,\; \mechanics,\; \Sigma,\; \mathcal{C} \bigr)
  \label{eq:task_tuple}
\end{equation}
%
where $\mathcal{S}$ is an HSSD scene~\citep{szot2021habitat}, $\agents = \{a_1, \ldots, a_n\}$ is a set of agents with spawn positions, $\goal$ is a \pddl{} goal formula combining physical predicates with $\K$-operators (Section~\ref{sec:epistemic_ops}), $\taskdesc$ is a natural-language task description shared with all agents, $\mechanics$ is a set of mechanic bindings drawn from the registry (Section~\ref{sec:mechanics}), $\Sigma$ maps each agent to private secrets, and $\mathcal{C} \in \{\text{cooperative}, \text{competitive}, \text{mixed}\}$ is the task category.

\paragraph{Secrets.} Each agent receives private secrets $\Sigma(a_i)$: natural-language statements of room restrictions, target object identities, and mechanic hints. Secrets state \emph{what} each agent privately knows but never \emph{how} to coordinate. Figuring out who to tell, what to say, and when to act is the epistemic coordination challenge the benchmark measures.

\paragraph{Categories.} \emph{Cooperative} tasks give all agents a shared $\goal$. \emph{Competitive} tasks assign two teams mutually exclusive win conditions over shared state. \emph{Mixed-motive} tasks combine a shared cooperative goal with private per-agent side-objectives that may conflict with the shared goal or with each other.

\subsection{Task generation agent}
\label{sec:gen_agent}

Each task is authored by a \textit{mini-SWE-style} \citep{minisweagent2025} coding agent. The agent operates inside a sandboxed workspace containing a blank task template (JSON file), the loaded scene graph (rooms, furniture, objects, agent spawns), reference files listing available mechanics, predicates, and actions, and a set of seed tasks sampled from the existing pool (see Appendix~\ref{app:workspace} for an example of workspace, the scene graph and prompt given to the coding agent).
The agent edits the task file using shell commands (\texttt{bash}, \texttt{jq}) and invokes five tool calls: \texttt{new\_scene}, \texttt{bash}, \texttt{judge}, \texttt{test\_task}, and \texttt{submit\_task}.

\paragraph{Inputs.} The agent receives four inputs at the start of each generation run: (a)~a sampled HSSD scene $\mathcal{S}$ with rooms, fixed furniture,  articulated furniture (openable/closable), movable objects, and agent spawn positions; (b)~a set of available mechanics (Section~\ref{sec:mechanics}); (c)~a target category $\mathcal{C} \in \{\text{cooperative}, \text{competitive}, \text{mixed}\}$ and target epistemic depth $d \in \{1,2,3\}$; and (d)~a set of seed tasks sampled from the existing pool with an 80/20 fail/pass ratio relative to the set of target models, biasing examples toward tasks that current frontier models cannot solve (Section~\ref{sec:calibration}).

\paragraph{Order of operations.} First, the agent writes the formal \pddl{} goal $\goal$ followed by the task description $\taskdesc$ and per-agent secrets $\Sigma(a_i)$ derived from $\goal$. This provides a formal structure to the task and prevents the narrative from containing requirements inconsistent with the final goal. Agent secrets consist of constraints and targets (room restrictions, object IDs, mechanic hints). We explicitly prompt the coding agent to never include the coordination strategy in secrets. The agents must figure out \emph{how} to communicate on their own. The judge (Section~\ref{sec:judge}) verifies that $\taskdesc$, $\Sigma$, and $\goal$ are consistent.

\paragraph{Iteration loop.} The coding agent invokes \texttt{judge} (Section~\ref{sec:judge}--\ref{sec:verification}) and \texttt{test\_task} (Section~\ref{sec:tom_necessity}) in any order and is free to iterate over the task by editing goals, adjusting mechanics, or calling \texttt{new\_scene} to start over. Each tool returns a structured feedback, the judge returns per-criterion scores and required fixes, the verifier returns planner diagnostics, and the calibrator returns pass/fail results . The agent uses this feedback to revise the task and resubmit. \texttt{submit\_task} requires both \texttt{judge} and \texttt{test\_task} to have passed. \g{ do we really want this line -- This closed-loop design distinguishes our approach from static pipelines that generate tasks first and filter afterward: the agent reasons about feedback and self-corrects before submission.}

\subsection{Quality judging}
\label{sec:judge}

A council of two LLMs (Kimi-K2.5 and GPT-5.2) independently scores each candidate task on eight criteria using a $[0, 1]$ rubric. \textit{Both models need to agree} for the task to pass. Seven criteria are shared across all categories:
\begin{itemize}[leftmargin=*]
  \item \textbf{Agent necessity:} every agent is central to the task, removing any one would break a required dependency or access path.
  \item \textbf{Secret quality:} secrets name exact target IDs and constraints but never prescribe coordination strategy.
  \item \textbf{Public/private grounding:} the shared task description $\taskdesc$ stays high-level; secrets $\Sigma$ carry the actionable specifics. This prevents broadcasting partial solutions to all agents.
  \item \textbf{Narrative consistency:} $\taskdesc$ and $\Sigma$ faithfully reflect the formal \pddl{} goal $\goal$.
  \item \textbf{Goal relevance:} every conjunct in $\goal$ contributes to the task; no filler predicates.
  \item \textbf{Mechanic utilization:} every binding in $\mechanics$ creates a distinct coordination challenge; no decorative mechanics.
  \item \textbf{Formal goal quality:} the \pddl{} specification is syntactically valid, solvable, and encodes meaningful epistemic structure.
\end{itemize}
An eighth, category-specific criterion is added. For cooperative tasks, \emph{task interdependence} requires that physical goals demand inter-agent information exchange. For competitive tasks, \emph{goal opposition} checks for mutually exclusive win conditions. For mixed tasks, \emph{subgoal tension} evaluates whether private objectives create genuine dilemmas. A task passes when:
%
\begin{equation}
  \bar{s} \geq \tau \quad \wedge \quad \forall\, c :\; s_c \geq 0.5
  \label{eq:judge_pass}
\end{equation}
%
where $\bar{s}$ is the mean score across all criteria, $\tau = 0.65$ is the overall threshold, and the per-criterion floor of $0.5$ prevents a single weak dimension from being masked by high scores elsewhere. Both thresholds were tuned on a held-out set of 50 manually rated tasks. On failure, the council returns per-criterion scores and concrete required fixes (e.g., ``add a secret for agent\_1 explaining the limited\_bandwidth mechanic''). The generation agent uses this feedback to revise and resubmit. The full judge prompt and an example of judge feedback are in Appendix~\ref{app:judge_prompt}. The rationale behind each criterion is discussed in Appendix~\ref{app:design_decisions}.

\g{later somewhere we want to add that what is the optimal strategy for mixed tasks and how it shows ToM
we also need to add in the paper everyehwere that the examples we show are completely random and not at all cherry picked}
\g{We need a human study on the outputs of the llm judge and we also need ablation of with and without judge tasks and then an ablation of the components of the judge. But, shd we wite this here on in experiments section? add}

\subsection{Formal verification \g{shd we name it formal verification?}}
\label{sec:verification}

The \texttt{`judge'} tool runs two deterministic checks. First, structural validation confirms that the \pddl{} goal is syntactically valid, that all referenced objects and furniture exist in the loaded scene, and that mechanic bindings are well-formed. Second, the $\K$-depth of the goal is computed by analyzing the nesting structure of epistemic operators in $\goal$ (Section~\ref{sec:epistemic_pddl}). If a target depth $d$ was specified for this generation run, the task is rejected if the computed depth does not match.%
\footnote{We encourage the reader to see Appendix~\ref{app:design_decisions} for why the planner is not given to the evaluated agent as a tool.}
\g{1. conduct a human study for the epistemic depth validity. 2. ablation of this judge}
% \g{the reviewer might have a question that why can't we directly use this solver as a tool call to llm, this is because this solver has entire information available whereas the llm needs to reason what is the information that other agent has. if it asks this question, this itself emans that the agent is trying to think about other agent's epistemic state, which means that it is trying to do tom reasoning. And if it does and asks the other agent about it's state then it has tom reasoning. So we just cannot give this solver to the llm as a tool. Don't write this here, write it in the design decision appendix, but put an * here for the reader and prompt them to look at the appendix.}

These checks are fast and deterministic. They can catch malformed tasks early, before the more expensive LLM evaluation and calibration runs. This verification step removes any task failure due to environmental failures, but physical executability (can agents actually navigate to objects, pick them up, and place them?) is not verified at this stage. That guarantee comes from the baseline calibration run (Section~\ref{sec:tom_necessity}).
% \g{we don't do this anymore}

\subsection{Empirical calibration}
\label{sec:tom_necessity}

\g{Structural?} verification confirms that a task is well-formed, but not that agents can physically execute it in the simulator. The empirical calibration step provides this guarantee. Each candidate task is run in a \emph{baseline} condition: all agent secrets are revealed to all agents, eliminating information asymmetry. If agents fail even with complete knowledge, the task either has a structural problem (like unreachable objects, impossible placements) or embodiment issues (cannot recognize the cup on the table) with the current agents. Only tasks that pass the baseline are admitted to the benchmark. This ensures that the task is solvable by the current agents and there is no environmental issues. 
% \g{this is done so that we can rule out any environment/embodied/solvability failures in the task to isolate the ToM failures. if an agent can complete the tas within the given constraints while knowing other agent's epistemic states then it measn that the task, in the actual eval, does not fail becuase of any environemntal issue, any embodiment challanges or rules out the possibility of task being not solvable. }

During evaluation, agents operate in \emph{standard} mode: each agent receives only its own secrets $\Sigma(a_i)$ and communicates through the bandwidth-limited channel. To account for the additional coordination overhead introduced by information asymmetry, we double the step budget and message limits relative to the baseline. One turn is for the agent to recognize what information asymmetry; one turn is to resolve that gap through communication. A task where the \textit{baseline} passes but the \text{standard} condition fails reflects the failure due to information asymmetry, i.e. 
% . : the physical task is solvable, the communication budget is sufficient, but 
agents do not model what their partners know.

\subsection{Difficulty evolution}
\label{sec:calibration}

A fixed benchmark saturates as models improve. \shortname{} addresses this by evolving the task pool over time through two mechanisms.

\paragraph{Seed-biased generation.} The generation agent receives seed tasks from the existing pool as in-context examples (Section~\ref{sec:gen_agent}). Seeds are sampled based on the target model's actual benchmark performance: 80\% of seeds are tasks the model failed, 20\% are tasks it passed. This biases the generator toward producing tasks that resemble the ones current models struggle with. As new tasks enter the pool and are benchmarked, the seed distribution shifts---later generation runs see harder examples, creating evolutionary pressure without changing the generation infrastructure.

\paragraph{Depth targeting.} Each generation run is assigned a target $\K$-depth $d \in \{1, 2, 3\}$, and submission rejects tasks whose computed depth does not match. Running generation in parallel across depth targets controls the epistemic depth profile of the benchmark.





% \saaket{Add short motivation}
% \section{Methodology}
% \label{sec:methodology}
% Task generation pipeline in \shortname{} has four stages: (1)~a ReAct-style LLM agent designs tasks through iterative scene exploration, (2)~a multi-model council judge scores candidates on structural and narrative criteria, (3)~a verifier (sylbolc + simulation based) confirms physical executability and ToM depth, and (4)~an evolution loop calibrates difficulty for the model being tested. We describe each stage below. Two worked task examples (cooperative and competitive) are in Appendix~\ref{app:examples}.

% \subsection{Task representation}
% \label{sec:task_repr}

% Each task is a self-contained specification. A task $\task$ is defined by:
% %
% \begin{equation}
%   \task = \bigl(\mathcal{S},\; \agents,\; \goal,\; \taskdesc,\; \mechanics,\; \Sigma,\; \mathcal{C} \bigr)
%   \label{eq:task_tuple}
% \end{equation}
% %
% where $\mathcal{S}$ is a Habitat Synthetic Scenes Dataset (HSSD) scene, $\agents = \{a_1, \ldots, a_n\}$ is a set of $n$ agents with spawn positions, $\goal$ is a \pddl{} goal formula, $\mechanics$ is a set of mechanic bindings, $\Sigma$ maps each agent to private secrets, and $\mathcal{C}$ is the task category. 

% The goal $\goal$ combines world-state predicates (\texttt{is\_on\_top}, \texttt{is\_inside}, \texttt{is\_open}, etc.) with the epistemic operator $\K$. $\K_i(\phi)$ asserts that agent $a_i$ \emph{knows} $\phi$, satisfied only when $a_i$ has directly observed $\phi$ or received it through communication. Goals nest $\K$ to depth $d \in \{0, 1, 2, 3\}$:
% %
% \begin{equation}
%   \underbrace{\K_{a_1}(\K_{a_2}(\cdots \K_{a_d}(\phi) \cdots))}_{\text{depth } d} \quad \Longrightarrow \quad \text{$(d{-}1)$-th order ToM}
%   \label{eq:tom_depth}
% \end{equation}

% We ground ToM depth in the level-$k$ reasoning framework from behavioral game theory~\citep{stahl1994experimental, camerer2004cognitive}. A level-0 agent follows a fixed heuristic without modeling others. A level-$k$ agent chooses the action that maximizes its expected utility assuming all other agents play level-$(k{-}1)$ policies:
% %
% \begin{equation}
%   \pi^{(0)}_i = \text{fixed heuristic}, \qquad \pi^{(k)}_i = \arg\max_{\pi_i} \; U_i\!\left(\pi_i,\; \pi^{(k-1)}_{-i}\right) \;\; \text{for } k \geq 1
%   \label{eq:level_k}
% \end{equation}
% %
% where $\pi^{(k-1)}_{-i}$ is the joint policy of all other agents at level $k{-}1$ and $U_i$ is agent $i$'s utility. The minimum level $k$ required for success is determined \emph{a priori} from the $\K$-nesting depth of $\goal$. A goal $\K_a(\K_b(\K_c(\phi)))$ requires level-3 reasoning: $a$ must model $b$'s model of $c$'s knowledge of $\phi$.

% \paragraph{Mechanics.} Mechanic bindings $\mechanics$ impose constraints that force epistemic reasoning. \textbf{Room restriction} confines agents to a subset of rooms. \textbf{Limited bandwidth} caps each agent's messages to $b_i$. \textbf{Restricted communication} specifies who can message whom via an adjacency graph. \textbf{Remote control} links a trigger object to a target, so acting on one changes the other's state in a different room. \textbf{State mirroring} synchronizes two objects' states. \textbf{Inverse state} reverses an object's handle (opening it closes it).

% \paragraph{Secrets.} Agent secrets $\Sigma(a_i)$ are natural-language strings that create information asymmetry: private knowledge, private objectives in mixed tasks, or team identities in competitive tasks.

% \paragraph{Task Category.} In \emph{cooperative} tasks, all agents share a single goal and succeed or fail together. In \emph{competitive} tasks, two teams have opposing goals over the same shared state (e.g., both teams want a contested object on their table). In \emph{mixed-motive} tasks, agents share a cooperative goal but each also has a private side-objective that it needs to achieve for task completion.

% \subsection{Task generation agent}
% \label{sec:gen_agent}

% The generator is a ReAct-style LLM agent that designs tasks through iterative reasoning, tool use, and self-correction. It follows a prescribed 9-step workflow: load a scene, inspect examples, draft the task JSON, verify PDDL correctness, submit for judge evaluation, iterate on feedback, verify the golden trajectory in the simulator, run an LLM-agent benchmark, and submit. Each step is gated: structural validity precedes qualitative review, which precedes empirical testing.

% The agent has access to six tools. \texttt{new\_scene} loads a random HSSD scene with $N$ agents ($\geq$5 locatable objects, $\geq$2 rooms). \texttt{bash} allows file exploration and task JSON editing within a sandboxed directory. \texttt{judge} first runs deterministic \pddl{} verification (syntax, object inventory, FastDownward solvability, ToM depth), then invokes the multi-model council (Section~\ref{sec:judge}). \texttt{verify\_golden\_trajectory} regenerates a deterministic action plan from the current spec and replays it in Habitat. \texttt{test\_task} runs a full LLM-agent benchmark and records calibration data. \texttt{submit\_task} requires all three gates---judge passed, trajectory verified, and benchmark run with nonzero progress---before writing the final task JSON.

% \subsection{Multi-model council judge}
% \label{sec:judge}

% Task quality is evaluated by a council of Kimi-K2.5 and GPT-5.2. Both models independently score the candidate on 8 criteria; both must agree for the task to pass. Disagreement defaults to rejection.

% The seven shared criteria are: \textbf{agent necessity} (removing any agent makes the task unsolvable), \textbf{secret quality} (actionable, no cross-agent leaks), \textbf{task naturalness} (no simulator IDs in descriptions), \textbf{narrative consistency} (description matches PDDL semantics), \textbf{goal relevance} (no extraneous predicates), \textbf{mechanic utilization} (every listed mechanic is required), and \textbf{PDDL solvability}. Each category adds one criterion: \emph{task interdependence} (cooperative), \emph{goal opposition} (competitive), or \emph{subgoal tension} (mixed). A task passes if, for both models $m$:
% %
% \begin{equation}
%   \bar{s}^{(m)} \geq \tau \quad \wedge \quad \forall\, i \neq i_{\text{nov}} :\; s_i^{(m)} \geq 0.5
%   \label{eq:judge_pass}
% \end{equation}
% %
% where $\bar{s}^{(m)}$ is the mean score, $\tau = 0.65$ (0.7 during evolution), and novelty $i_{\text{nov}}$ is exempt from the per-criterion floor.

% \subsection{Verification}
% \label{sec:sim_verify}

% Every task passes three checks before acceptance.


% \begin{minipage}[t]{0.46\linewidth}
% \paragraph{PDDL verification.} The verifier parses the goal, compiles it against the scene's object inventory, and runs FastDownward (with PDKB fallback for epistemic predicates) to confirm solvability. It computes the ToM depth by counting $\K$-nesting. For the goal in Figure~\ref{fig:pddl_example}, the verifier extracts $\K$-depth 2, checks that \texttt{bowl\_1}, \texttt{table\_22}, and \texttt{cabinet\_34} exist in the scene, and confirms a valid action sequence achieves all conjuncts.
% \end{minipage}
% \hfill
% \begin{minipage}[t]{0.5\linewidth}
% \begin{lstlisting}
% (and
%   (is_on_top bowl_1 table_22)
%   (K agent_0 (K agent_1
%     (is_on_top bowl_1 table_22)))
%   (is_open cabinet_34)
% )
% \end{lstlisting}
% \vspace{-0.5em}
% \captionof{figure}{Example PDDL goal with epistemic nesting. The $\K$-depth is 2 (first-order ToM).}
% \label{fig:pddl_example}
% \end{minipage}%

% \paragraph{Golden trajectory verification.} The verifier generates a deterministic action plan from the PDDL goal (Navigate, Pick, Place, Open, Close, Communicate) and replays it step-by-step in Habitat. \g{Here we are also adding another check that if full information is given agents must be able to solve the tasks}. The final simulator state is checked against $\goal$. Trajectories are regenerated each time rather than cached, so any change to the task specification is immediately tested.

% \paragraph{Specification validation.} Structural checks run without the simulator: template placeholder detection, mechanic schema validation, and room-restriction consistency (no agent's golden trajectory navigates to a restricted room).

% \begin{figure}[h]
% \begin{center}
% \resizebox{0.95\linewidth}{!}{\input{figures/evolution_pipeline_tikz}}
% \end{center}
% \caption{Evolution pipeline. The system benchmarks the task pool against a model ladder (GPT-5-mini $\to$ Sonnet $\to$ GPT-5.2), generates harder tasks when $\rho > \rho^*$, and repeats until the pool is calibrated.}
% \label{fig:evolution}
% \end{figure}

% \subsection{Evolution pipeline}
% \label{sec:evolution}

% The evolution pipeline calibrates task difficulty against a model ladder ($\text{GPT-5-mini} \to \text{Sonnet} \to \text{GPT-5.2}$). At each tier, the pool-wide pass rate $\rho$ is computed. If $\rho > \rho^*$ (default 0.20), the pool is too easy: an in-context learning sampler selects representative failures (9 failed + 1 passed task, sorted by completion percentage) as few-shot examples and generates harder tasks. The sampler adjusts its guidance based on $\rho$: below 0.10, it shifts to a different difficulty axis; between 0.10 and 0.30, it amplifies failure patterns; above 0.30, it targets model weaknesses directly. Generation halts when $\rho \leq \rho^*$ or the budget is exhausted.

% The pipeline targets a distribution of approximately 30\% level-1, 45\% level-2, and 25\% level-3 tasks (8\% tolerance), concentrating at level-2 where we observe the best signal for discriminating model capabilities.

% \subsection{Visual observation pipeline}
% \label{sec:vision}
% \saaket{this belong in experimental setup for describing the baseline}
% Each agent perceives the environment through an egocentric RGB camera in Habitat. At every turn, the simulator renders a first-person video of the agent's surroundings. Rather than feeding the full video into the LLM context, we use a two-stage key-frame selection process: video frames are subsampled to 10--15 candidates, then the agent's LLM selects 2--5 frames most informative for its current reasoning step given its full context (instructions, communication history, prior observations). The selected frames enter the LLM context as images for the next action decision.

% Vision creates a tighter coupling between perception and ToM. An agent that sees a teammate pick up an object has direct evidence; an agent restricted from a room gets no visual information about it and must rely on communicated beliefs. This asymmetry is what makes epistemic goals non-trivial: the $\K$-nesting depth measures how many relay hops are needed to propagate a visual observation to an agent who cannot see it.

\section{Agent Behavioral Analysis}
\label{sec:behaviors}

\begin{table*}[t]
\centering
\small
\renewcommand{\arraystretch}{1.25}
\setlength{\tabcolsep}{4.5pt}
\begin{tabular}{@{} l cc c cc c cc @{}}
\toprule
& \multicolumn{2}{c}{\textbf{Cooperative}} & \phantom{a}
& \multicolumn{2}{c}{\textbf{Mixed-Motive}} & \phantom{a}
& \multicolumn{2}{c}{\textbf{Overall}} \\
\cmidrule{2-3} \cmidrule{5-6} \cmidrule{8-9}
\textbf{Model}
& \textbf{Functional} & \textbf{Literal}
&& \textbf{Functional} & \textbf{Literal}
&& \textbf{Functional} & \textbf{Literal} \\
\midrule
Claude 4.5 Sonnet \\
Claude 4.5 Haiku            & 39.3 & 36.2 && \textbf{52.0} & 24.8 && \textbf{45.7} & 30.6 \\
Claude 4.5 Opus             & \textbf{44.0} & \textbf{45.6} && 49.2$^\dagger$ & \textbf{41.7}$^\dagger$ && 45.5$^\dagger$ & \textbf{44.5}$^\dagger$ \\
\hline
GPT-5.4          & 13.3 & --   && 26.7 & --   && 20.0 & --   \\
GPT-5.4-mini     & 17.6 & 40.5 && 22.1 & 24.4 && 19.9 & 32.5 \\
\hline
Gemini-3.1-Pro\\
Gemini-3-Flash\\
\hline
DeepSeek-v3.2    & 38.7 & 38.1 && 35.1 & 34.2 && 36.9 & 36.2 \\
Kimi-K2.5        & 33.3 & 47.5 && 38.7 & 38.9 && 36.0 & 43.2 \\
O3               & 28.0 & \textbf{55.6} && 24.0 & 35.0 && 26.0 & 45.4 \\

\bottomrule
\end{tabular}

\vspace{2pt}
\caption{Functional vs.\ literal Theory of Mind across frontier models. Functional ToM measures whether agents succeed at tasks requiring epistemic reasoning under partial observability. Literal ToM measures whether agents can answer explicit questions about other agents' beliefs during the same tasks. O3 and Kimi-K2.5 score highest on literal ToM but rank below Haiku and Opus on functional ToM, indicating that reporting beliefs and acting on them are distinct capabilities.}
\label{tab:main_results}
\end{table*}

We evaluate GPT-5.2 (vision mode, 35~tasks) and Sonnet (text mode, 33~tasks) on EmToM-generated tasks spanning 0th--3rd order epistemic depth, 2--4~agents, and cooperative, mixed-motive, and competitive settings. GPT-5.2 passes 6/35 tasks (17\%); Sonnet passes 8/33 (24\%). Rather than reporting aggregate accuracy, we characterize seven behavioral patterns that define the current capability frontier. Full task descriptions, agent instructions, and verbatim dialogue traces for each finding are in Appendix~\ref{app:traces}.

\subsection{Second-order ToM: achievable but fragile}
\label{sec:behav_tom2}

GPT-5.2 achieves genuine second-order belief reasoning---$\K(a_i, \K(a_j, \varphi))$---in 4 of 35~tasks. These successes share a common profile: ring communication topology, 2--4~agents, and $\geq$4 messages per agent. In the \emph{Manifest Chain} task (Appendix~\ref{app:manifest}), $a_0$ must end knowing that $a_3$ knows the cushion remains on its chair. Agent $a_0$ plans the full second-order relay in a single message, $a_2$ verifies the fact and explicitly grants $a_3$ justified belief (``you can truthfully say you know this''), and $a_3$ asserts knowledge with provenance (``I KNOW \ldots agent\_2 verified''). The agents distinguish \emph{receiving a message} from \emph{holding verified knowledge}.

This capability is fragile. At depth~3 or bandwidth $\leq$2 messages per agent, no model succeeds. Sonnet shows 0~instances of second-order ToM across all 33~tasks; its belief modeling stays at the action level (``they will move X'') rather than the epistemic level (``they think we don't know X'').

\subsection{Cooperative bias: zero defection}
\label{sec:behav_coop_bias}

Across 12 mixed-motive tasks in both model runs, no agent pursued a private objective that conflicted with the shared goal---a 100\% cooperation rate. In \emph{Silent Alarm Calibration} (Appendix~\ref{app:silent_alarm}), each of three agents receives a private side-request that directly contradicts one shared sub-goal. All three agents follow the shared goal without deliberation; the private objectives never appear in reasoning traces.

The cooperative bias extends beyond passive compliance. In 4 of 12 mixed-motive tasks, agents voluntarily \emph{disclose} their private objectives. Agent $a_0$ broadcasts ``(my side objective)'' verbatim; $a_1$ writes ``personal objective---please do NOT move picture\_frame\_2'' (Appendix~\ref{app:transparency}). This is a compound ToM failure: agents neither pursue their own interest nor model that revealing private information changes others' strategic reasoning. The zero-defection pattern holds across both GPT-5.2 and Sonnet, suggesting it is model-agnostic---likely an artifact of RLHF training that optimizes for cooperation.

\subsection{Confirmation spirals}
\label{sec:behav_spiral}

In relay tasks, agents fall into recursive confirmation loops where every agent requests verification from the next before acting. In \emph{Privacy Handoff Relay} (Appendix~\ref{app:spiral}), only $a_3$ can enter the bathroom. Instead of scouting, $a_3$ asks $a_2$ for the information that $a_2$ just asked $a_3$ to provide. All 8~messages are consumed on meta-communication; zero ground-truth observations are relayed. Agent $a_3$ finally scouts at turn~31 but cannot relay findings---all message budgets are exhausted.

The root cause is a role-awareness failure: agents cannot distinguish ``I need information from someone'' from ``I am the source of information for others.'' A related pattern appears in \emph{Triple-Consensus Protocol} (Appendix~\ref{app:deadlock}), where two agents mutually wait for the other's confirmation after both have exhausted their message budgets.

\subsection{Physical--epistemic decoupling}
\label{sec:behav_decouple}

Agents treat epistemic relay goals and physical manipulation goals as independent sequential tasks. In \emph{Relay-of-Belief Drill} (Appendix~\ref{app:decoupling}), a three-hop belief relay executes perfectly by turn~22, but only 1~of~3 required cabinets is opened. Agents allocate all communication bandwidth to the relay and have no messages left for physical task coordination. 65 of 81~turns are \texttt{Wait} actions.

At depth~1, 80\% of failures are physical (wrong object, wrong receptacle). At depth~2+, the pattern inverts: 70\%+ of failures are epistemic (broken relay, exhausted budget). The evolution pipeline exploits this: tightening bandwidth pushes tasks from physically to epistemically challenging without changing the physical layout.

\subsection{Competitive counter-play}
\label{sec:behav_competitive}

Competitive 2v2 tasks produce four emergent strategy types (Appendix~\ref{app:competitive}): (i)~\emph{object theft}---agents pick up contested items from the opponent's goal location; (ii)~\emph{cabinet guarding}---an agent stations itself at a shared resource to re-open it whenever the opponent closes it; (iii)~\emph{item concealment}---hiding contested objects inside closed containers; and (iv)~\emph{remote-control exploitation}---toggling linked objects across rooms.

All four strategies are reactive---agents respond to the opponent's most recent action. No agent plans multiple moves ahead or models the opponent's likely counter-strategy. In \emph{Claim the Glass} (Appendix~\ref{app:stalemate}), teams steal each other's placed objects in alternating turns, producing a stalemate where neither team completes the goal.

\subsection{Action repetition loops}
\label{sec:behav_loops}

Sonnet exhibits a failure mode absent in GPT-5.2: indefinite repetition of failed actions. In \emph{Mirror-Fridge Stance War} (Appendix~\ref{app:loops}), $a_1$ sends the identical message 22~consecutive times (turns 9--30) despite receiving ``you have used all 2 of your allowed messages'' after each attempt. In \emph{Audit Code} (Appendix~\ref{app:loops}), $a_0$ tries \texttt{Navigate[cabinet\_96]} 15~times despite being told the room is off-limits---a constraint stated in its own instructions.

GPT-5.2 shows a milder form (max 4~repeats), suggesting better feedback integration. A related Sonnet failure is tool confusion: in \emph{Locker Rally}, agents use \texttt{FindObjectTool} (for movable objects) when the correct tool is \texttt{FindReceptacleTool} (for furniture), and never try the alternative.

\subsection{Model comparison}
\label{sec:behav_comparison}

\begin{table}[h]
\centering
\small
\begin{tabular}{lcc}
\toprule
\textbf{Dimension} & \textbf{GPT-5.2 (Vision)} & \textbf{Sonnet (Text)} \\
\midrule
Pass rate & 6/35 (17\%) & 8/33 (24\%) \\
Avg.\ turns (pass) & $\sim$11 & $\sim$7 \\
Wait actions & 40\% & 16\% \\
Message repetition & max 4 & 19--22 \\
Room restriction loops & none & up to 15 \\
2nd-order ToM & 4 instances & 0 instances \\
Defection rate & 0\% & 0\% \\
\bottomrule
\end{tabular}
\caption{Behavioral comparison of GPT-5.2 and Sonnet on EmToM tasks.}
\label{tab:model_comparison}
\end{table}

GPT-5.2 shows deeper epistemic reasoning (4~instances of second-order ToM) and better feedback integration (max 4~action repeats) but is significantly more passive (40\% \texttt{Wait}). Sonnet is faster when it succeeds ($\sim$7 vs.\ $\sim$11 turns) but degenerates catastrophically under failure---22 identical messages, 15 blocked navigations. Neither model handles tasks with 3+ mechanics or 4+ agents reliably. The communication budget is the primary bottleneck: agents spend messages on coordination overhead rather than task-critical information.

%\section{Evaluation}
%\section{Results}
%\section{Conclusion}

\bibliography{colm2026_conference}
\bibliographystyle{colm2026_conference}
\newpage
\appendix
\section{Orders of Theory of Mind}
\label{app:tom_orders}

\definecolor{tomboxbg}{RGB}{245, 247, 252}
\definecolor{tomboxframe}{RGB}{180, 190, 215}

\begin{tcolorbox}[
  enhanced,
  colback=tomboxbg,
  colframe=tomboxframe,
  boxrule=0.5pt,
  arc=3pt,
  left=4pt, right=4pt, top=4pt, bottom=4pt,
  title={\small\bfseries Orders of Theory of Mind},
  coltitle=black,
  colbacktitle=tomboxbg,
  toptitle=2pt, bottomtitle=2pt,
  fonttitle=\sffamily,
  attach boxed title to top left={yshift=-2mm, xshift=4mm},
  boxed title style={colback=tomboxbg, colframe=tomboxframe, boxrule=0.5pt, arc=2pt},
  float=t,
]
\footnotesize
\renewcommand{\arraystretch}{1.4}
\setlength{\tabcolsep}{4pt}
\begin{tabularx}{\linewidth}{@{} >{\raggedright\arraybackslash}p{1.1cm} >{\centering\arraybackslash}p{2.4cm} >{\raggedright\arraybackslash}X >{\raggedright\arraybackslash}X @{}}

\textbf{Order} & \textbf{Formal} & \textbf{Everyday analogy} & \textbf{EmToM example} \\[2pt]
\midrule\\[-6pt]

\textbf{0} \newline \textit{\scriptsize No ToM}
& $\phi$
& You open the fridge because you are thirsty. No one else is involved.
& Agent navigates to a cabinet and opens it. No coordination needed. \\[4pt]

\textbf{1} \newline \textit{\scriptsize Self-aware}
& $\mathcal{K}_a(\phi)$
& You can't reach the top shelf, so you ask your roommate to check if the jar is there.
& Agent~0 is barred from the kitchen. It asks Agent~1 to check the fridge and report back. It models \emph{its own} gap, not Agent~1's mind. \\[4pt]

\textbf{2} \newline \textit{\scriptsize Other-aware}
& $\mathcal{K}_a\!\bigl(\mathcal{K}_b(\phi)\bigr)$
& Your roommate left before the package arrived. You know \emph{they} don't know it's here, so you text them.
& Agent~0 must ensure Agent~2 knows the bowl is on the table. Agent~2 hasn't seen it. Agent~0 must reason about what Agent~2 lacks and relay it. \\[4pt]

\textbf{3} \newline \textit{\scriptsize Recursive}
& $\mathcal{K}_a\!\bigl(\mathcal{K}_b\!\bigl(\mathcal{K}_c(\phi)\bigr)\bigr)$
& You need your manager to know the client got the contract. Your assistant sent it, so you ask them to confirm they told the manager.
& Agent~0 must know that Agent~1 knows that Agent~2 knows the wardrobe is open. Requires a planned relay chain across three agents. \\[4pt]

\textbf{4+} \newline \textit{\scriptsize Self-ref.}
& $\mathcal{K}_a\!\bigl(\mathcal{K}_b\!\bigl(\cdots\bigr)\bigr)$
& Poker: you think your opponent thinks you're bluffing, so you bet high---but they may know that you know they think this.
& Goal: $\mathcal{K}_{a_1}\!\bigl(\mathcal{K}_{a_0}\!\bigl(\mathcal{K}_{a_1}\!\bigl(\mathcal{K}_{a_0}(\phi)\bigr)\bigr)\bigr)$---a self-referential loop through two agents. \\

\end{tabularx}

\vspace{4pt}
\noindent\rule{\linewidth}{0.3pt}
\vspace{2pt}
{\scriptsize\textit{Frontier models reliably handle order~1 but fail at order~2 and above (Section~\ref{sec:behaviors}).}}
\label{tab:tom_orders}
\end{tcolorbox}
\section{Agent Trajectory Evidence}
\label{app:traces}

This appendix provides full task descriptions, agent instructions, and verbatim dialogue traces for each behavioral finding reported in Section~\ref{sec:behaviors}. All quotes are extracted directly from agent \texttt{thought}, \texttt{action}, and \texttt{result} fields in the benchmark logs.

\subsection{Manifest Chain: Second-Order ToM}
\label{app:manifest}

\paragraph{Task.} \emph{Manifest Chain: Second-Order Confirmation + Cross-Room Switchboard.} 4~agents, directed ring ($a_0 \to a_1 \to a_2 \to a_3 \to a_0$), 2~messages per agent. Mechanics: \texttt{restricted\_communication}, \texttt{limited\_bandwidth}, \texttt{room\_restriction}. Result: \textbf{PASSED} in 17~turns.

\paragraph{Agent instructions.}
\begin{itemize}
\item $a_0$: Cannot enter kitchen, living room, or bedroom~3. Can only message $a_1$. Must finish with second-order certainty: needs a message from $a_3$ explicitly stating they \emph{know} the living-room cushion is still on its chair.
\item $a_1$: Cannot enter recreation room or living room. Can only message $a_2$. Must verify picture frame is still on bedroom-3 table. Must receive a message establishing that they know $a_2$ knows the picture-frame fact.
\item $a_2$: Cannot enter bedroom~3. Can only message $a_3$. Only agent who can enter living room: must verify cushion is still on chair. Must move the kitchen cup to the recreation-room table.
\item $a_3$: Cannot enter recreation room, living room, or bedroom~3. Can only message $a_0$. Must send $a_0$ an explicit statement that they know the cushion is on its chair.
\end{itemize}

\paragraph{Key turns.}

\begin{quote}\small
\textbf{T1} $a_0 \to a_1$: ``Please coordinate confirmations: (A) have agent\_2 verify the living-room cushion is still on its chair, then tell agent\_3; then get a message back from agent\_3 stating they KNOW the cushion is still on its chair (I need that). (B) have someone verify the picture frame is still on bedroom-3 table \ldots''
\end{quote}

\begin{quote}\small
\textbf{T3} $a_2$: \texttt{FindObjectTool[cushion on the living room chair]} $\to$ \texttt{cushion\_4 on chair\_38 in living\_room\_1 2.54m away}
\end{quote}

\begin{quote}\small
\textbf{T4} $a_2 \to a_3$: ``Verified in living\_room\_1: cushion\_4 is still on chair\_38 (I did not move it). You can truthfully say you know this.''
\end{quote}

\begin{quote}\small
\textbf{T5} $a_3 \to a_0$: ``I KNOW the living-room cushion is still on the living-room chair---agent\_2 verified in living\_room\_1 that cushion\_4 is still on chair\_38 and did not move it.''
\end{quote}

The relay chain completes by turn~5. Turns 8--17 are spent on the physical sub-task (cup relocation), during which three agents idle with \texttt{Wait} while $a_2$ searches for entity names.


\subsection{Silent Alarm: Zero Defection}
\label{app:silent_alarm}

\paragraph{Task.} \emph{Silent Alarm Calibration with Private Side-Requests.} 3~agents. Mechanics: \texttt{room\_restriction}, \texttt{limited\_bandwidth}. Result: \textbf{PASSED} in 5~turns.

\paragraph{Agent instructions and conflicts.} Shared goal: hallway stand OPEN, living-room table CLOSED, kitchen cabinet OPEN. Each agent has a conflicting private objective:
\begin{itemize}
\item $a_0$: private---wants living-room table to end OPEN (shared goal: CLOSED).
\item $a_1$: private---wants hallway stand to end CLOSED (shared goal: OPEN).
\item $a_2$: private---wants kitchen cabinet to end CLOSED (shared goal: OPEN).
\end{itemize}

\paragraph{Key turns.}

\begin{quote}\small
\textbf{T1} $a_0$ (broadcast): ``Agent\_1: Please open the tall kitchen cabinet \ldots Agent\_2: Please locate hallway display stand and open it; confirm when done.'' \\
\emph{No mention of private objective in thought or message.}
\end{quote}

\begin{quote}\small
\textbf{T4} $a_2$: \texttt{Open[stand\_34]} $\to$ ``Successful execution!'' \\
\emph{$a_2$'s private goal was to close this stand. Ignored without deliberation.}
\end{quote}

\begin{quote}\small
\textbf{T5} $a_1$: \texttt{Open[cabinet\_39]} $\to$ ``Successful execution!'' \\
\emph{$a_1$'s private goal was to close this cabinet. Ignored without deliberation.}
\end{quote}

All three agents follow the shared task. No private objective appears in any reasoning trace.


\subsection{Strategy Transparency: Private Objective Disclosure}
\label{app:transparency}

\paragraph{Task.} \emph{Relay-Signature Configuration with Hidden Side Aims.} 4~agents. Mechanics: \texttt{restricted\_communication}, \texttt{limited\_bandwidth}, \texttt{room\_restriction}. Result: \textbf{FAILED} (21~turns).

\paragraph{Agent instructions.} $a_0$: cannot enter kitchen or bedroom~2; private objective is to keep the kitchen refrigerator OPEN. $a_3$: can enter bedroom~2; private objective is to place a second lamp on the bed.

\paragraph{Key turn: private objective disclosure.}

\begin{quote}\small
\textbf{T1} $a_0 \to a_1$: ``I can't enter kitchen or bedroom2. Please relay to the agent in bedroom2: confirm the stable starting observation there \ldots Also, please have someone with kitchen access make sure the refrigerator ends up OPEN \textbf{(my side objective)}.''
\end{quote}

\begin{quote}\small
\textbf{T2} $a_0 \to a_1$: [identical message repeated verbatim, consuming second and final message]
\end{quote}

$a_0$ labels its private objective with the parenthetical ``(my side objective)'' in a message visible to teammates. It then sends the same message twice, wasting its entire bandwidth on a single duplicated request. Meanwhile, $a_3$ quietly completes its non-conflicting private objective (placing the lamp on the bed) without disclosure.


\subsection{Confirmation Spiral: Privacy Handoff Relay}
\label{app:spiral}

\paragraph{Task.} \emph{Privacy Handoff Relay: Move the Stuffed Toy and Verify the Bathroom Cabinet (ToM-3).} 4~agents. Mechanics: \texttt{room\_restriction}, \texttt{restricted\_communication}, \texttt{limited\_bandwidth}. Result: \textbf{FAILED} (21~turns).

\paragraph{Agent instructions.} $a_0$: cannot enter bathroom or office; holds the stuffed toy. $a_1$: relay node. $a_2$: relay node; cannot enter bathroom or office. $a_3$: \textbf{only agent} who can enter bathroom and office; designated scout. Each agent has 2~messages.

\paragraph{The spiral.} All four agents use their first message to request confirmation from the next agent in the chain:

\begin{quote}\small
\textbf{T1} $a_0 \to a_1$: ``Need relay-confirmed scout report: has someone (agent\_2/3) visually checked BOTH (1)~bathroom cabinet status (must be left open) and (2)~office shelves/key items?''

\textbf{T1} $a_1 \to a_2$: ``Need confirmation relay: please verify you SEE (1) stuffed toy and its exact object name/location \ldots''

\textbf{T1} $a_2 \to a_3$: ``Need exact confirmations for relay: did you personally see (1) the stuffed toy \ldots (3) the bathroom cabinet?''

\textbf{T1} $a_3 \to a_2$: ``Before I act: confirm you personally saw BOTH key items---(1) the stuffed toy (where exactly?) and (2) the bathroom cabinet (is its door open or closed?)?''
\end{quote}

$a_3$---the only agent with bathroom access---asks $a_2$ for the information $a_2$ just asked $a_3$ to provide. All 8 messages consumed. Zero ground-truth observations relayed.

\paragraph{Post-spiral deadlock.} $a_3$ opens cabinet\_48 at turn~11 but has no messages to report it. $a_0$ holds the stuffed toy for 16~turns but never places it, waiting for a relay-confirmed report that never comes. Turns 16--21: all agents execute \texttt{Wait[None]}.


\subsection{Consensus Deadlock: Triple-Consensus Protocol}
\label{app:deadlock}

\paragraph{Task.} \emph{Triple-Consensus Protocol: Verify Two Remote States, Then Seal the Report.} 3~agents, one-way ring ($a_0 \to a_1 \to a_2 \to a_0$), 2~messages each. Mechanics: \texttt{restricted\_communication}, \texttt{limited\_bandwidth}, \texttt{room\_restriction}. Result: \textbf{FAILED} (21~turns).

\paragraph{Physical tasks.} Completed correctly---$a_2$ opens the spray-bottle cabinet and closes the bowl-container.

\paragraph{Epistemic deadlock.} $a_2$ reports the bowl-container state to $a_0$ (turn~9). $a_0$ forwards it to $a_1$ (turn~10), consuming both messages. Then $a_2$ reports the spray-bottle cabinet state to $a_0$ (turn~10)---but $a_0$ has 0~messages left and cannot relay ``I know'' to $a_1$.

\begin{quote}\small
\textbf{T21} $a_0$ (Done): ``I still need agent\_1's explicit acknowledgment that they know I know the spray-bottle cabinet is open \ldots''

\textbf{T21} $a_1$ (Done): ``I still need agent\_0 to explicitly tell me `I KNOW the spray-bottle cabinet is open,' and I can't send any more messages \ldots''
\end{quote}

Both agents identify the deadlock in their final reasoning but cannot resolve it---a mutual wait on messages that can never be sent.


\subsection{Physical--Epistemic Decoupling: Relay-of-Belief Drill}
\label{app:decoupling}

\paragraph{Task.} \emph{Relay-of-Belief Drill: Three Cabinets and a Verified Hallway Fact.} 4~agents. Mechanics: \texttt{restricted\_communication}, \texttt{room\_restriction}, \texttt{limited\_bandwidth}. Result: \textbf{FAILED} (20~turns).

\paragraph{Agent instructions.} $a_3$ can enter the hallway; must observe the bowl's location and relay through the chain. Physical goal: open 3~cabinets (kitchen, bathroom~1, bathroom~2). $a_0$: 2~messages. $a_1$: 3~messages. $a_2$: 3~messages. $a_3$: 2~messages.

\paragraph{Epistemic relay (perfect).}
\begin{quote}\small
\textbf{T4} $a_3 \to a_2$: ``Hallway fact (verbatim): bowl\_1 is ON table\_22 in hallway\_1.''

\textbf{T5} $a_2 \to a_1$: ``Hallway fact (verbatim from agent\_3): bowl\_1 is ON table\_22 in hallway\_1.''

\textbf{T6} $a_1 \to a_0$: ``Relay-of-belief: agent\_2 confirmed (from agent\_3, verbatim) the hallway fact: `bowl\_1 is ON table\_22 in hallway\_1.'\,''
\end{quote}

Three-hop relay with explicit provenance tracking, completed by turn~6.

\paragraph{Physical tasks (failed).} Only \texttt{cabinet\_34} (kitchen) opened by $a_2$ at turn~4. Neither bathroom cabinet is ever touched. No agent coordinates who can access the bathrooms. 14 of 20~turns (7--20) are \texttt{Wait[None]} across all agents.


\subsection{Competitive Counter-Play and Stalemates}
\label{app:competitive}

\paragraph{Object theft: Signal vs.\ Staging.} 4~agents (2v2). Mechanics: \texttt{restricted\_communication}, \texttt{limited\_bandwidth}, \texttt{room\_restriction}. Result: \textbf{FAILED} (20~turns). GPT-5.2 vision run.

\begin{quote}\small
\textbf{T3} $a_3$ (Team~1): \texttt{Pick[bottle\_4]} from \texttt{chest\_of\_drawers\_38}

\textbf{T5} $a_1$ (Team~0): ``In bedroom\_2, the special bottle is bottle\_4 and it is currently on agent\_3 (opposing) about 1.2m from me.''

\textbf{T5} $a_0 \to a_1$: ``\ldots bottle\_4 on chest\_of\_drawers\_38 in bedroom\_2 (likely the special bottle). I can't navigate there---can you retrieve bottle\_4 and bring it to a living room table \ldots''

\textbf{T7} $a_3$: \texttt{Place[bottle\_4, on, bed\_31]} --- Team~1's staging spot.

\textbf{T11} $a_1$: \texttt{Pick[bottle\_4]} from \texttt{bed\_31} --- \emph{steals} from opponent's staging spot.

\textbf{T13} $a_1$: \texttt{Place[bottle\_4, on, table\_11]} --- Team~0's staging spot.
\end{quote}

Team~0 detects the opponent holding the contested asset, issues a competitive directive, and executes a theft. Both teams then idle (turns~14--20). Task fails because neither team satisfies the storage-unit constraint.


\paragraph{Cabinet guarding: Display Stance Bowl Duel.} 4~agents (2v2). Sonnet run. Result: \textbf{FAILED} (20~turns).

Team~0 needs \texttt{cabinet\_34} CLOSED; Team~1 needs it OPEN.
\begin{quote}\small
\textbf{T5} $a_2$ (Team~1): \texttt{Open[cabinet\_34]} $\to$ success.

\textbf{T9} $a_0$ (Team~0): \texttt{Close[cabinet\_34]} $\to$ success.

\textbf{T10} $a_2$ (Team~1): \texttt{Open[cabinet\_34]} $\to$ success. \\
$a_2$ then executes \texttt{Wait[None]} for turns~11--20, stationed at the cabinet to guard it.
\end{quote}

\paragraph{Item concealment: Switchboard Standoff.} 4~agents (2v2). Sonnet run. Result: \textbf{FAILED} (20~turns).

\begin{quote}\small
\textbf{T7} $a_2$ (Team~1): \texttt{Pick[tray\_2]}

\textbf{T10} $a_2$: \texttt{Open[cabinet\_31]}

\textbf{T11} $a_2$: \texttt{Place[tray\_2, within, cabinet\_31]} --- hides tray inside cabinet.

\textbf{T12} $a_2$: \texttt{Close[cabinet\_31]} --- conceals the tray.
\end{quote}

\label{app:stalemate}
\paragraph{Object race stalemate: Claim the Glass.} 4~agents (2v2). Sonnet run. Result: \textbf{FAILED} (20~turns).

Team~0 needs glass + plate on dining table. Team~1 needs both on TV-room table.
\begin{quote}\small
\textbf{T5} $a_0$: \texttt{Pick[glass\_6]} --- Team~0 secures glass.

\textbf{T9} $a_3$: \texttt{Pick[plate\_0]} --- Team~1 secures plate.

\textbf{T15} $a_1$: \texttt{Pick[plate\_0]} from TV table --- Team~0 steals plate.

\textbf{T17} $a_3$: \texttt{Pick[glass\_6]} from dining table --- Team~1 steals glass.
\end{quote}

Final state: glass on Team~1's table, plate on Team~0's table. Neither team has both items. Stalemate.


\subsection{Action Repetition Loops}
\label{app:loops}

\paragraph{Message loop: Mirror-Fridge Stance War.} 4~agents (2v2). Sonnet run. Result: \textbf{FAILED} (30~turns). $a_1$ has a 2-message budget.

\begin{quote}\small
\textbf{T8} $a_1 \to a_0$: ``I need to handle bottle placement on smaller living room table \ldots''

\textbf{T9} $a_1 \to a_0$: ``Team\_0 target: bottle on SMALLER living room table (table\_12 or table\_19), fridge OPEN, living room counter\_36 OPEN. I'll handle bottle and counter.''

\textbf{T10--T30} $a_1$: [identical \texttt{Communicate} action repeated 21~times] $\to$ ``You have used all 2 of your allowed messages. You cannot send any more.''
\end{quote}

$a_1$ never executes a physical action. 21~consecutive turns wasted on the same failed \texttt{Communicate} call.

\paragraph{Room restriction loop: Audit Code.} 4~agents (2v2). Sonnet run. Result: \textbf{FAILED} (20~turns). $a_0$ cannot enter \texttt{bathroom\_1}.

\begin{quote}\small
\textbf{T5--T16} $a_0$: \texttt{Navigate[cabinet\_96]} $\to$ ``You cannot enter bathroom\_1 to reach cabinet\_96. The area is off-limits to you.'' \\
{[12 consecutive identical attempts]}
\end{quote}

\paragraph{Tool confusion: Locker Rally.} 4~agents (2v2). Sonnet run. Result: \textbf{FAILED} (20~turns).

$a_3$ (Team~1) searches for wardrobes using the wrong tool:
\begin{quote}\small
\textbf{T5} $a_3$: \texttt{FindObjectTool[wardrobes in laundryroom\_1]} $\to$ ``No objects found''

\textbf{T6} $a_3$: \texttt{FindObjectTool[wardrobe in laundryroom\_1]} $\to$ ``No objects found''

\textbf{T10} $a_3$: \texttt{FindObjectTool[wardrobe or armoire or dresser]} $\to$ ``No objects found''

\emph{\ldots 10+ failed attempts. Never tries \texttt{FindReceptacleTool}.}
\end{quote}

Meanwhile, $a_1$ (Team~0) succeeds on the first try:
\begin{quote}\small
\textbf{T6} $a_1$: \texttt{FindReceptacleTool[wardrobe]} $\to$ ``wardrobe\_29 in laundryroom, wardrobe\_30 in laundryroom''
\end{quote}

Wardrobes are classified as receptacles, not objects. The simulator distinguishes movable objects (\texttt{FindObjectTool}) from static furniture (\texttt{FindReceptacleTool}). Sonnet agents frequently select the wrong tool and never recover.


\section{Task Examples}
\label{app:examples}

\subsection{Cooperative: depth-4 epistemic chain}
\label{app:example_coop}

Four agents coordinate a ``silent inspection'' across a household: open the kitchen cabinet, close the bedroom wardrobe, and place items in designated locations. Communication follows a linear chain ($a_0 \leftrightarrow a_1 \leftrightarrow a_2 \leftrightarrow a_3$) with one message per agent. Room restrictions prevent any single agent from accessing all locations; only $a_1$ can enter the kitchen.

The PDDL goal includes:
%
\begin{equation}
  \K_{a_0}\bigl(\K_{a_1}\bigl(\K_{a_2}\bigl(\K_{a_3}(\texttt{is\_on\_top}(\texttt{box}, \texttt{cabinet}))\bigr)\bigr)\bigr)
  \label{eq:example_k4}
\end{equation}
%
This requires $a_0$ to know that $a_1$ knows that $a_2$ knows that $a_3$ knows the box is on the cabinet. The box is in the kitchen, visible only to $a_1$. For this chain to hold, $a_1$ must relay the observation through $a_2$ to $a_3$, consuming two of the chain's four message slots. Agent $a_0$ must infer that $a_1$ relayed rightward (the only viable direction), that $a_2$ forwarded it, and that $a_1$ anticipated this relay. This is third-order ToM.

\begin{table}[h]
\begin{center}
\small
\begin{tabular}{lllp{4.2cm}}
\toprule
\textbf{Agent} & \textbf{Can message} & \textbf{Restricted from} & \textbf{Key private knowledge} \\
\midrule
$a_0$ & $a_1$ & kitchen & Cabinet holds a cardboard box \\
$a_1$ & $a_0$ or $a_2$ & entryway, bedroom, living room & Only agent with kitchen access \\
$a_2$ & $a_1$ or $a_3$ & entryway, bathroom, kitchen & Must relay between $a_1$ and $a_3$ \\
$a_3$ & $a_2$ & kitchen & Can access bathroom (target surface) \\
\bottomrule
\end{tabular}
\end{center}
\caption{Agent configuration for the cooperative silent inspection task. Each agent has 1 message. The linear chain and room restrictions force depth-4 epistemic reasoning.}
\label{tab:example_agents}
\end{table}

\subsection{Competitive: self-referential beliefs}
\label{app:example_comp}

Two teams ($\{a_0, a_1\}$ vs.\ $\{a_2, a_3\}$) compete to set a home's ``final inspection'' configuration. Each team must place a contested cushion in their preferred location and set a kitchen latch to their required state. Communication is team-internal only, with 1--2 messages per agent.

The winning condition for team~0 includes:
%
\begin{equation}
  \K_{a_1}\bigl(\K_{a_0}\bigl(\K_{a_1}\bigl(\K_{a_0}(\texttt{is\_closed}(\texttt{latch}))\bigr)\bigr)\bigr)
  \label{eq:example_selfref}
\end{equation}
%
This is a self-referential depth-4 chain: $a_1$ must know that $a_0$ knows that $a_1$ knows that $a_0$ knows the latch is closed. The belief chain loops through the same two agents, requiring $a_1$ to model $a_0$'s model of $a_1$'s own knowledge.

Room restrictions add complexity: $a_0$ and $a_2$ cannot enter the kitchen, while $a_1$ and $a_3$ cannot enter rooms containing the remote-control trigger and inverse-state stand. Agents must reason about teammate beliefs and about what the opposing team believes about their strategy.

\section{Epistemic compilation example}
\label{app:compilation}

We walk through the full compilation of the epistemic goal from Figure~\ref{fig:pddl_example}. The task has two agents (\texttt{agent\_0} and \texttt{agent\_1}), where \texttt{agent\_1} is in the same room as \texttt{bowl\_1} and \texttt{table\_22}, and \texttt{agent\_1} can communicate with \texttt{agent\_0} with a message budget of 2.

\paragraph{Original goal (with epistemic operators).}

\begin{lstlisting}
(and
  (is_on_top bowl_1 table_22)
  (K agent_0 (K agent_1
    (is_on_top bowl_1 table_22)))
  (is_open cabinet_34)
)
\end{lstlisting}

The planner needs to verify that (1) the bowl can be placed on the table, (2) the cabinet can be opened, and (3) \texttt{agent\_0} can come to know that \texttt{agent\_1} knows the bowl has been placed. The physical subgoals (1) and (2) are already classical \pddl{}. The epistemic subgoal (3) is not. The compiler transforms it as follows.

\paragraph{Step 1: Create knowledge predicates.} For the leaf fact \texttt{(is\_on\_top bowl\_1 table\_22)}, the compiler creates a first-layer knowledge predicate for each agent:

\begin{lstlisting}
(knows_agent_0_a1b2c3d4)   ; agent_0 knows bowl is on table
(knows_agent_1_a1b2c3d4)   ; agent_1 knows bowl is on table
\end{lstlisting}

Because the goal nests two $\K$-operators, the compiler also creates a second-layer predicate for the outer agent:

\begin{lstlisting}
(knows_agent_0_e5f6g7h8)   ; agent_0 knows that agent_1
                            ; knows bowl is on table
\end{lstlisting}

\paragraph{Step 2: Create observe operators.} \texttt{Agent\_1} is co-located with the bowl and table, so the compiler generates an observe operator that sets \texttt{agent\_1}'s knowledge predicate when the physical fact holds:

\begin{lstlisting}
(:action observe_knows_agent_1_a1b2c3d4
  :parameters ()
  :precondition (is_on_top bowl_1 table_22)
  :effect (knows_agent_1_a1b2c3d4))
\end{lstlisting}

This operator says: if the bowl is on the table and \texttt{agent\_1} is in the room (encoded in the precondition via the observability model), then \texttt{agent\_1} can be marked as knowing this fact. No real observation happens; the planner is simply checking whether a valid sequence of such operators exists.

\paragraph{Step 3: Create inform operators.} \texttt{Agent\_1} can communicate with \texttt{agent\_0}, so the compiler generates inform operators that propagate first-layer knowledge. Each operator consumes one message token:

\begin{lstlisting}
(:action inform_knows_agent_0_a1b2c3d4_from_agent_1_tok1
  :parameters ()
  :precondition (and
    (knows_agent_1_a1b2c3d4)
    (can_communicate agent_1 agent_0)
    (msg_tok_agent_1_1))
  :effect (and
    (knows_agent_0_a1b2c3d4)
    (not (msg_tok_agent_1_1))))
\end{lstlisting}

A second copy uses \texttt{msg\_tok\_agent\_1\_2}, giving the planner two opportunities (matching the budget of 2). The \texttt{(not (msg\_tok\_...))} effect consumes the token, preventing reuse.

\paragraph{Step 4: Create nested-knowledge inform operators.} For the outer $\K$-goal, the compiler generates an operator that lets \texttt{agent\_1} inform \texttt{agent\_0} about its own knowledge state:

\begin{lstlisting}
(:action inform_knows_agent_0_e5f6g7h8_from_agent_1_tok2
  :parameters ()
  :precondition (and
    (knows_agent_1_a1b2c3d4)
    (can_communicate agent_1 agent_0)
    (msg_tok_agent_1_2))
  :effect (and
    (knows_agent_0_e5f6g7h8)
    (not (msg_tok_agent_1_2))))
\end{lstlisting}

The precondition requires that \texttt{agent\_1} already knows the fact (first layer). The effect sets the second-layer predicate, establishing that \texttt{agent\_0} now knows \texttt{agent\_1} knows.

\paragraph{Step 5: Replace the goal.} The original epistemic goal is replaced with a conjunction of classical predicates:

\begin{lstlisting}
(and
  (is_on_top bowl_1 table_22)
  (knows_agent_1_a1b2c3d4)
  (knows_agent_0_e5f6g7h8)
  (is_open cabinet_34)
)
\end{lstlisting}

\paragraph{Step 6: Add budget tokens to the initial state.} The problem's \texttt{:init} section is augmented with tokens representing \texttt{agent\_1}'s message budget:

\begin{lstlisting}
(msg_tok_agent_1_1)
(msg_tok_agent_1_2)
\end{lstlisting}

\paragraph{Result.} The compiled problem is entirely classical \pddl{}. Fast Downward searches over the physical actions (place, open, navigate) together with the observe and inform operators. A valid plan might be:

\begin{enumerate}
  \item \texttt{agent\_1} places \texttt{bowl\_1} on \texttt{table\_22}.
  \item \texttt{observe\_knows\_agent\_1\_a1b2c3d4} fires (agent\_1 sees the bowl is placed).
  \item \texttt{inform\_knows\_agent\_0\_a1b2c3d4\_from\_agent\_1\_tok1} fires (agent\_1 tells agent\_0 the fact, consuming token 1).
  \item \texttt{inform\_knows\_agent\_0\_e5f6g7h8\_from\_agent\_1\_tok2} fires (agent\_1 tells agent\_0 that it knows, consuming token 2).
  \item Another agent opens \texttt{cabinet\_34}.
\end{enumerate}

If Fast Downward finds this plan (or any valid alternative), the task is provably solvable. The $\K$-depth of 2 is read directly from the nesting structure during Step~1.

\section{Task generation agent workspace and prompt}
\label{app:workspace}

The generation agent operates in an isolated workspace directory:
\begin{lstlisting}
workspace/
  working_task.json      # task being authored
  template.json          # blank task skeleton
  current_scene.json     # rooms, furniture, objects, spawns
  sampled_tasks/         # seed tasks (80% fail, 20% pass)
  submitted_tasks/       # accepted tasks
  available_mechanics.md # mechanic registry
  available_predicates.md
  available_actions.md
\end{lstlisting}

\paragraph{Example scene graph.} The \texttt{current\_scene.json} file describes the loaded Habitat scene. Below is an abridged example for a 2-agent, 5-room scene:

\begin{lstlisting}
{
  "scene_id": "102344280",
  "episode_id": "885",
  "rooms": ["office_1", "dining_room_1",
            "laundryroom_1", "kitchen_2",
            "entryway_1"],
  "furniture": ["table_25", "cabinet_33",
       "couch_9", "counter_30", "cabinet_31",
       "table_18"],
  "objects": ["cushion_2", "bowl_4"],
  "articulated_furniture": ["cabinet_33",
                            "cabinet_31"],
  "furniture_in_rooms": {
    "office_1":      ["table_25", "cabinet_33"],
    "dining_room_1": ["couch_9"],
    "laundryroom_1": ["counter_30"],
    "kitchen_2":     ["cabinet_31"],
    "entryway_1":    ["table_18"]
  },
  "objects_on_furniture": {
    "table_25": ["cushion_2"],
    "chair_10": ["bowl_4"]
  },
  "agent_spawns": {
    "agent_0": {"position": [1.2, 0.1, 3.4],
                "room": "office_1"},
    "agent_1": {"position": [5.6, 0.1, 2.1],
                "room": "dining_room_1"}
  }
}
\end{lstlisting}

The agent uses this to select goal-relevant objects and furniture, configure room restrictions, and ground the \pddl{} goal in actual scene entities. Only articulated furniture (cabinets, fridges, drawers) can appear in \texttt{is\_open}/\texttt{is\_closed} goals.

\paragraph{System prompt.} The agent receives the following system prompt (abridged; full prompt is ${\sim}$400 lines):

\definecolor{promptbg}{RGB}{248, 246, 240}
\definecolor{promptframe}{RGB}{200, 195, 175}

\begin{tcolorbox}[
  enhanced,
  colback=promptbg,
  colframe=promptframe,
  boxrule=0.5pt,
  arc=3pt,
  left=5pt, right=5pt, top=5pt, bottom=5pt,
  title={\small\bfseries Generation Agent System Prompt (abridged)},
  coltitle=black,
  colbacktitle=promptbg,
  toptitle=2pt, bottomtitle=2pt,
  fonttitle=\sffamily,
  attach boxed title to top left={yshift=-2mm, xshift=4mm},
  boxed title style={colback=promptbg, colframe=promptframe, boxrule=0.5pt, arc=2pt},
  breakable,
]
\scriptsize
\ttfamily

You are a puzzle designer creating multi-agent collaboration challenges.\\[4pt]

\textnormal{\sffamily\bfseries Response format.}
Each turn: \texttt{Thought: [reasoning]} then \texttt{Action: tool\_name[argument]}.\\[4pt]

\textnormal{\sffamily\bfseries Tools.}\\
\texttt{new\_scene[N]} -- load HSSD scene with N agents, reset task.\\
\texttt{bash[cmd]} -- run shell commands (jq, cat, python3, \ldots).\\
\texttt{judge[]} -- PDDL verification + golden trajectory replay + LLM quality evaluation.\\
\texttt{test\_task[]} -- run standard + baseline calibration.\\
\texttt{submit\_task[]} -- save task (requires judge + test\_task).\\[4pt]

\textnormal{\sffamily\bfseries Workflow.}\\
1. \texttt{new\_scene[N]} $\to$ load scene.\\
2. Inspect seed tasks in \texttt{sampled\_tasks/} for inspiration.\\
3. Edit \texttt{working\_task.json}: author the \texttt{problem\_pddl :goal} FIRST, then write \texttt{task}, \texttt{agent\_secrets}, and mechanic bindings to match. Do NOT hand-author \texttt{:objects} or \texttt{:init}.\\
4. \texttt{judge[]} $\to$ fix $\to$ repeat until pass.\\
5. \texttt{test\_task[]} $\to$ calibrate difficulty.\\
6. \texttt{submit\_task[]}.\\[4pt]

\textnormal{\sffamily\bfseries Core rules.}\\
-- Author the PDDL goal as the source of truth; write narrative to match it.\\
-- Secrets state WHAT (room restrictions, target IDs, mechanic hints) but NEVER HOW (no coordination strategy, no relay instructions).\\
-- Every agent must make a distinct, non-substitutable contribution.\\
-- At least one physical action must be information-dependent: an agent cannot determine what to do without information held by another agent.\\
-- Do not prescribe coordination strategy in secrets. The agent must figure out how to communicate.\\[4pt]

\textnormal{\sffamily\bfseries Secret examples.}\\[2pt]
\textnormal{\sffamily BAD} (leaks coordination strategy):\\
\texttt{agent\_0: "Wait for agent\_3 to tell you whether stand\_34 is open, then forward that to agent\_0."}\\[2pt]
\textnormal{\sffamily GOOD} (states constraints, agents figure out coordination):\\
\texttt{agent\_0: ["You cannot enter hallway\_2.", "You can only message agent\_1. You can send 2 messages.", "By the end, you must be confident a teammate knows stand\_34 is open."]}\\[4pt]

\textnormal{\sffamily\bfseries Functional ToM patterns} (use at least one):\\
1. \textnormal{Delegation choice} -- agent must choose which teammate to inform; only one can act on the info.\\
2. \textnormal{Sequencing choice} -- correct action order depends on what a teammate already knows.\\
3. \textnormal{Relay choice} -- sender cannot reach the actor directly; must pick a relay path.\\
4. \textnormal{Information-gated action} -- agent's correct action depends on a fact only another agent can observe.\\
5. \textnormal{Mixed-motive cooperation} -- private objectives change how useful or reliable a teammate is.\\
6. \textnormal{Competitive blocking} -- best play depends on inferring the opponent's likely branch.\\[4pt]

\textnormal{\sffamily\bfseries Self-tests before submitting.}\\
-- Remove all Communicate actions from the golden trajectory. Can agents still achieve all physical goals independently? If yes, the task does not test functional ToM.\\
-- Can the outermost K() agent directly walk to the room and observe the fact? If yes, the K-goal is trivial.\\
\end{tcolorbox}

\section{Judge council prompt and feedback}
\label{app:judge_prompt}

The judge council (Kimi-K2.5 and GPT-5.2) receives the following prompt (abridged):

\begin{tcolorbox}[
  enhanced,
  colback=promptbg,
  colframe=promptframe,
  boxrule=0.5pt,
  arc=3pt,
  left=5pt, right=5pt, top=5pt, bottom=5pt,
  title={\small\bfseries Judge Council Prompt (abridged)},
  coltitle=black,
  colbacktitle=promptbg,
  toptitle=2pt, bottomtitle=2pt,
  fonttitle=\sffamily,
  attach boxed title to top left={yshift=-2mm, xshift=4mm},
  boxed title style={colback=promptbg, colframe=promptframe, boxrule=0.5pt, arc=2pt},
  breakable,
]
\scriptsize
\ttfamily
You are an expert evaluator for multi-agent tasks.\\[4pt]

\textnormal{\sffamily\bfseries Context provided to the judge:}\\
-- Task category and category-specific rules\\
-- Available actions, mechanics, items, and predicates\\
-- Scene objects (rooms, furniture, objects on furniture)\\
-- The compiled formal view (PDDL after mechanic compilation)\\
-- Derived runtime semantics (functional goal + literal-ToM probes)\\
-- The full task JSON (description, secrets, mechanics, PDDL goal)\\[4pt]

\textnormal{\sffamily\bfseries Key checks:}\\
-- \texttt{task} is global and must not leak secret targets\\
-- Secrets must be actionable (exact IDs) and must not prescribe strategy\\
-- Every K() goal must be backed by a mechanic preventing direct observation\\
-- The functional projection (after dropping K-goals) must remain non-trivial\\
-- K() probes should test who knows functionally relevant facts under real asymmetry\\
-- Reward tasks where the best action depends on a partner-specific model; penalize pure fact-relay tasks\\[4pt]

\textnormal{\sffamily\bfseries Scoring:} Each criterion scored 0.0--1.0.\\[2pt]
\textnormal{\sffamily\bfseries Output:} JSON with per-criterion scores, reasoning (under 15 words each), and required fixes.\\
\end{tcolorbox}

\paragraph{Example judge feedback.} Below is an actual judge response for a task that failed on secret quality and mechanic utilization (task selected at random from the generation logs):

\begin{tcolorbox}[
  enhanced,
  colback=promptbg,
  colframe=promptframe,
  boxrule=0.5pt,
  arc=3pt,
  left=5pt, right=5pt, top=5pt, bottom=5pt,
  breakable,
]
\scriptsize
\ttfamily
\{\\
\quad "agent\_necessity": \{"score": 0.7, "reasoning": "Both agents have distinct room access"\},\\
\quad "secret\_quality": \{"score": 0.3, "reasoning": "Agent\_0 secret prescribes relay strategy"\},\\
\quad "task\_naturalness": \{"score": 0.8, "reasoning": "Clean high-level description"\},\\
\quad "narrative\_consistency": \{"score": 0.7, "reasoning": "Matches PDDL goal"\},\\
\quad "goal\_relevance": \{"score": 0.8, "reasoning": "All conjuncts needed"\},\\
\quad "mechanic\_utilization": \{"score": 0.4, "reasoning": "room\_restriction is decorative"\},\\
\quad "pddl\_solvability": \{"score": 0.9, "reasoning": "Valid and solvable at K-2"\},\\
\quad "task\_interdependence": \{"score": 0.6, "reasoning": "Some parallel execution possible"\},\\
\quad "overall\_reasoning": "Secret leaks strategy; room restriction does not block any goal.",\\
\quad "required\_fixes": [\\
\qquad "Remove strategy hints from agent\_0 secret",\\
\qquad "Restrict agent\_0 from kitchen\_1 where cabinet\_28 is"\\
\quad ]\\
\}
\end{tcolorbox}

\section{Key design decisions}
\label{app:design_decisions}

Several design choices in the generation pipeline were informed by failure modes observed during early development. We document each decision and the problem it addresses.

\paragraph{Why agent necessity.} Early generated tasks often included agents that contributed nothing---an agent would be assigned to a room but have no goal-relevant object there, or two agents would have identical access and capabilities, making one redundant. When benchmarked, the redundant agent would simply idle (Wait actions for the entire episode) while the other completed the task alone. The \emph{agent necessity} criterion rejects tasks where any agent can be removed without breaking the intended solution. This forces the generator to design tasks where each agent holds unique access, information, or physical capability.

\paragraph{Why secrets must not prescribe strategy.} In early experiments, secrets contained instructions like ``ask agent\_1 to open the fridge and relay the result to agent\_2.'' Agents followed these instructions verbatim, achieving near-perfect pass rates---but the task reduced to instruction following, not epistemic reasoning. The agent never needed to model what others know or can do; the secret told it exactly what to communicate and to whom. The \emph{secret quality} criterion now rejects any secret that leaks coordination strategy. Secrets state constraints (``you cannot enter kitchen\_1''), targets (``cabinet\_28 must end open''), and mechanic hints (``the handle is reversed'')---never the plan.

\paragraph{Why public/private grounding.} If the shared task description $\taskdesc$ contains exact object IDs and target locations, all agents receive the same complete information and there is no reason to communicate. Information asymmetry arises only when $\taskdesc$ stays high-level (``reset the house for inspection'') and the actionable specifics (which cabinet, which room, which object) are distributed across private secrets $\Sigma(a_i)$. This split is what creates the need for agents to share information selectively.

\paragraph{Why PDDL goal before narrative.} When the generation agent wrote the natural-language description first, it frequently invented requirements not expressible in \pddl{} (e.g., ``agents should feel satisfied with the arrangement'') or omitted requirements that were in the formal goal. Writing $\goal$ first and deriving $\taskdesc$ and $\Sigma$ from it eliminated this class of inconsistencies.

\paragraph{Why a two-model council.} A single judge model exhibited systematic biases: GPT-5.2 was lenient on secret quality (rarely flagging strategy leakage), while Kimi-K2.5 was lenient on mechanic utilization (accepting decorative mechanics). Requiring both models to agree compensates for each model's blind spots. Tasks that pass the council satisfy a stricter quality bar than either model alone would enforce.

\paragraph{Why baseline calibration.} The baseline condition (all secrets public) serves two purposes. First, it proves the task is physically solvable---if agents fail even with full information, the task has a structural problem (unreachable objects, impossible goal states). Second, the gap between baseline (pass) and standard (fail) isolates the contribution of information asymmetry. Without this control, we cannot distinguish ``the task is hard because it requires epistemic coordination'' from ``the task is hard because the objects are hard to find.''

\paragraph{Why not give the planner to the evaluated agent.} The epistemic planner (Section~\ref{sec:verification}) has access to all agents' secrets, room restrictions, and the complete goal formula simultaneously. It solves the task as an omniscient coordinator. If we gave this planner as a tool to the evaluated agent, the agent could query it to determine what every other agent knows, what information to communicate, and in what order---bypassing epistemic reasoning entirely. The whole point of the benchmark is that the agent must \emph{infer} what its partners know from their room access, communication history, and behavior. Handing the agent an oracle that answers ``what does agent\_1 know?'' would reduce the task to tool calling, not Theory of Mind.


\end{document}
